% =====================================================================
% DSO-style Sensitivity analysis for the Loyola (LO) index.
% Honest / publication-safe. Uses only pre-specified reductions + pure-gamma
% points. Figures: fig_lo_prediction_correlation_heatmap, fig_lo_intercorrelation_heatmap,
% fig_lo_degeneracy_order10_trees, fig_lo_structure_sensitivity_decanes.
% Table: \begin{table}[H]
\centering
\caption{Pearson correlation $r$ between selected Loyola reductions and pre-specified pure-$\gamma$ Loyola points and the five physicochemical properties of the 18 octane isomers (NIST WebBook). Bold = best $|r|$ per property. Edge count $m=LO(0,0,0)$ is omitted (constant on octanes).}
\label{tab:lo_octane}
\small
\begin{tabular}{lrrrrr}
\toprule
Index & $T_B$ & $\Delta H_f$ & $\Delta H_{\rm vap}$ & $S$ & $\omega$\\
\midrule
$M_1{=}LO(0,1,0)$ & -0.718 & -0.762 & -0.912 & -0.973 & -0.970\\
$M_2{=}LO(1,0,0)$ & -0.503 & -0.553 & -0.775 & -0.942 & \textbf{-0.987}\\
$HM{=}LO(0,2,0)$ & -0.655 & -0.712 & -0.873 & \textbf{-0.975} & -0.982\\
${}^mM_2{=}LO(-1,0,0)$ & \textbf{+0.861} & \textbf{+0.888} & +0.942 & +0.866 & +0.799\\
$R{=}LO(-1/2,0,0)$ & +0.821 & +0.848 & +0.956 & +0.935 & +0.901\\
$\chi{=}LO(0,-1/2,0)$ & +0.801 & +0.829 & +0.953 & +0.950 & +0.927\\
$H/2{=}LO(0,-1,0)$ & +0.822 & +0.846 & +0.959 & +0.935 & +0.904\\
$ISI{=}LO(1,-1,0)$ & -0.037 & -0.047 & -0.363 & -0.632 & -0.760\\
$GA{=}2\,LO(1/2,-1,0)$ & +0.823 & +0.856 & +0.959 & +0.942 & +0.909\\
$AG{=}\frac{1}{2}LO(-1/2,1,0)$ & -0.817 & -0.856 & -0.951 & -0.941 & -0.903\\
$LO(0,0,1)$ & -0.828 & -0.831 & -0.966 & -0.931 & -0.930\\
$LO(0,0,2)$ & -0.830 & -0.844 & \textbf{-0.967} & -0.939 & -0.926\\
\bottomrule
\end{tabular}
\end{table}

% =====================================================================
\section{Sensitivity analysis}\label{sec:sensitivity}

Following the structure of the diminished-Sombor study
\cite{MovahediGutmanRedzepovicFurtula2026}, we benchmark the Loyola index against
the classical degree-based indices it reduces to and report prediction ability,
intercorrelation, degeneracy, and structure sensitivity. To keep the analysis
honest and reproducible, we use \emph{only} pre-specified parameter values: the
reductions $M_1,M_2,HM,{}^mM_2,R,\chi,H/2,ISI,GA,AG$ (Section~\ref{sec:intro}) and
the two pre-specified pure-$\gamma$ points $LO(0,0,1)$ and $LO(0,0,2)$. No
parameter is tuned after inspecting the property correlations; the reasons for
excluding such decimal tuning are given in the supplementary note. The edge count
$m=LO(0,0,0)$ is omitted throughout because it is constant on the octanes.

\subsection{Prediction ability}\label{sub:prediction}
Table~\ref{tab:lo_octane} and Figure~\ref{fig:lo_pred} report the Pearson
correlation between each index and five physicochemical properties of the 18
octane isomers (boiling point $T_B$, enthalpy of formation $\Delta H_f$, heat of
vaporization $\Delta H_{\mathrm{vap}}$, entropy $S$, acentric factor $\omega$;
properties from the NIST WebBook). The best $|r|$ per property is highlighted. The
modified second Zagreb index attains the strongest correlation on $T_B$ and
$\Delta H_f$, the hyper-Zagreb index on $S$, and the second Zagreb index on
$\omega$; the pre-specified pure-$\gamma$ point $LO(0,0,2)$ attains the strongest
correlation on $\Delta H_{\mathrm{vap}}$ $(|r|=0.967)$. The Loyola family is thus
competitive and, at selected pre-specified parameter settings, improves the
correlation for some properties. (The properties here are from the NIST WebBook;
\cite{MovahediGutmanRedzepovicFurtula2026} used a different source, so the
$\Delta H_{\mathrm{vap}}$ and $S$ correlations may differ by data source rather
than by index computation.)

\begin{table}[H]
\centering
\caption{Pearson correlation $r$ between selected Loyola reductions and pre-specified pure-$\gamma$ Loyola points and the five physicochemical properties of the 18 octane isomers (NIST WebBook). Bold = best $|r|$ per property. Edge count $m=LO(0,0,0)$ is omitted (constant on octanes).}
\label{tab:lo_octane}
\small
\begin{tabular}{lrrrrr}
\toprule
Index & $T_B$ & $\Delta H_f$ & $\Delta H_{\rm vap}$ & $S$ & $\omega$\\
\midrule
$M_1{=}LO(0,1,0)$ & -0.718 & -0.762 & -0.912 & -0.973 & -0.970\\
$M_2{=}LO(1,0,0)$ & -0.503 & -0.553 & -0.775 & -0.942 & \textbf{-0.987}\\
$HM{=}LO(0,2,0)$ & -0.655 & -0.712 & -0.873 & \textbf{-0.975} & -0.982\\
${}^mM_2{=}LO(-1,0,0)$ & \textbf{+0.861} & \textbf{+0.888} & +0.942 & +0.866 & +0.799\\
$R{=}LO(-1/2,0,0)$ & +0.821 & +0.848 & +0.956 & +0.935 & +0.901\\
$\chi{=}LO(0,-1/2,0)$ & +0.801 & +0.829 & +0.953 & +0.950 & +0.927\\
$H/2{=}LO(0,-1,0)$ & +0.822 & +0.846 & +0.959 & +0.935 & +0.904\\
$ISI{=}LO(1,-1,0)$ & -0.037 & -0.047 & -0.363 & -0.632 & -0.760\\
$GA{=}2\,LO(1/2,-1,0)$ & +0.823 & +0.856 & +0.959 & +0.942 & +0.909\\
$AG{=}\frac{1}{2}LO(-1/2,1,0)$ & -0.817 & -0.856 & -0.951 & -0.941 & -0.903\\
$LO(0,0,1)$ & -0.828 & -0.831 & -0.966 & -0.931 & -0.930\\
$LO(0,0,2)$ & -0.830 & -0.844 & \textbf{-0.967} & -0.939 & -0.926\\
\bottomrule
\end{tabular}
\end{table}


\begin{figure}[H]\centering
\includegraphics[width=0.8\linewidth]{fig_lo_prediction_correlation_heatmap.pdf}
\caption{Correlation of selected Loyola reductions and pre-specified pure-$\gamma$ points with octane properties. Best absolute correlation per property is marked ($\ast$).}
\label{fig:lo_pred}\end{figure}

\subsection{Intercorrelation}\label{sub:intercorr}
Figure~\ref{fig:lo_inter} shows the intercorrelation matrix among the selected
indices on the 18 octanes. On this small, degree-limited set the pure-$\gamma$
points are highly collinear with the classical reductions: the maximum absolute
correlation of $LO(0,0,1)$ and $LO(0,0,2)$ with a classical reduction is $0.983$
and $0.993$, respectively. This indicates that the $\gamma$-imbalance term does
\emph{not} add independent linear information on the octane benchmark.

\begin{figure}[H]\centering
\includegraphics[width=0.7\linewidth]{fig_lo_intercorrelation_heatmap.pdf}
\caption{Intercorrelation among selected Loyola reductions and pure-$\gamma$ points on the octane isomer set. The pure-$\gamma$ points are highly collinear with the classical reductions.}
\label{fig:lo_inter}\end{figure}

\subsection{Degeneracy}\label{sub:degeneracy}
Degeneracy measures an index's ability to discriminate among non-isomorphic
graphs; lower is better. On the $106$ non-isomorphic trees of order $10$ we use
$\mathrm{Degeneracy}(\%) = 100\,(1 - \#\text{distinct}/N)$ (Figure~\ref{fig:lo_deg}).
The pure-$\gamma$ points $LO(0,0,1)$ and $LO(0,0,2)$ have degeneracy $25.5\%$,
among the lowest of the compared indices and comparable to $R$, $GA$, and $AG$;
$M_1$ ($83.0\%$) and $M_2$ ($67.0\%$) are the most degenerate. (The exact
degeneracy convention used by \cite{MovahediGutmanRedzepovicFurtula2026} is to be
verified against its source before final insertion.)

\begin{figure}[H]\centering
\includegraphics[width=0.85\linewidth]{fig_lo_degeneracy_order10_trees.pdf}
\caption{Degeneracy percentage of selected indices on the $106$ order-$10$ trees. Lower degeneracy indicates better discrimination.}
\label{fig:lo_deg}\end{figure}

\subsection{Structure sensitivity}\label{sub:structure}
We compute structure sensitivity $SS=\sigma/\mu$, abruptness
$Abr=(\max-\min)/\mu$, and $SA=SS/Abr$ for the index values on the $75$ decane
isomers (Figure~\ref{fig:lo_ss}). The pure-$\gamma$ point $LO(0,0,2)$ has the
highest structure sensitivity ($SS=0.165$) and abruptness among the compared
indices, with $LO(0,0,1)$ also high; the $\gamma$-imbalance term therefore
contributes most visibly to structure sensitivity rather than to added
correlation information. (These are standard structure-sensitivity formulas; they
do not reproduce the numerical values of
\cite{MovahediGutmanRedzepovicFurtula2026} and are to be reconciled with the
Raki\'c--Furtula method before final insertion.)

\begin{figure}[H]\centering
\includegraphics[width=0.95\linewidth]{fig_lo_structure_sensitivity_decanes.pdf}
\caption{Structure-sensitivity metrics on the $75$ decane isomers. Formulas for $SS$, $Abr$, and $SA$ are stated in the methods note and should be verified against \cite{MovahediGutmanRedzepovicFurtula2026} before final insertion.}
\label{fig:lo_ss}\end{figure}

\paragraph{Summary.} The Loyola family unifies several classical degree-based
indices as exact reductions. In the octane benchmark, selected pre-specified
Loyola values were competitive with the classical reductions (best on
$\Delta H_{\mathrm{vap}}$), while the pure-$\gamma$ values showed high structure
sensitivity but limited independent correlation information on this small,
degree-limited set.
